% source: J. S. Milne, "Abelian Varieties" (course notes), v2.00, 2008, www.jmilne.org
% pdf_page: 0109
% doc_page: 103 (Chapter III, Section 5 — end of proofs of Lemma 5.4, Proposition 5.3 and Theorem 5.1; Corollary 5.5, Remark 5.6)
% retrieved: 2026-07-16
% transcribed_by: claude-fable-5 (vision, rendered at 200dpi via pdftoppm)

% [running head: 5. CANONICAL MAPS, page 103]

\DeclareMathOperator{\Ima}{Im}    % printed as "Im"
\DeclareMathOperator{\Ker}{Ker}   % printed as "Ker"
\DeclareMathOperator{\Div}{Div}   % Div^D_C, the divisor functor

% [continuation of the proof of Lemma 5.4 from the previous page]

and on multiplying this by $(-1)^n X_i^n$ and summing over $n$ (so that the
successive terms cancel out) we obtain the identity
\[
  \sigma_m - \sigma_{m-1} X_i + \dots + (-1)^m X_i^m = \sigma_m(i).
\]
On multiplying this with $dX_i$ and summing, we get the required identity.
\hfill $\square$

We now complete the proof of (5.3). First let $D = rQ$. Then
$\hat{\mathcal{O}}_Q = k[[X]]$ and
$\hat{\mathcal{O}}_D = k[[\sigma_1, \dots, \sigma_r]]$ (see the proof of
(3.2); by $\mathcal{O}_D$ we mean the local ring at the point $D$ on
$C^{(r)}$). If $\omega = (a_0 + a_1 X + a_2 X^2 + \cdots) dX$, $a_i \in k$,
when regarded as an element of $\Omega^1_{\hat{\mathcal{O}}_Q / k}$, then
$\omega' = a_0 \tau_0 + a_1 \tau_1 + \cdots$. We know that
$\{d\sigma_1, \dots, d\sigma_r\}$ is a basis for
$\Omega^1_{\hat{\mathcal{O}}_D / k}$ as an $\hat{\mathcal{O}}_D$-module, but
the lemma shows that $\tau_0, \dots, \tau_{r-1}$ is also a basis. Now
$(\omega) \ge D$ and $\omega'(D) = 0$ are both obviously equivalent to
$a_0 = a_1 = \dots = a_{r-1} = 0$. The proof for other divisors is similar.

We finally prove the exactness of the sequence in (5.1). The injectivity of
$(di)_D$ follows from the fact that $i \colon F \hookrightarrow C^{(r)}$ is a
closed immersion. Moreover the sequence is a complex because $f \circ i$ is
the constant map $x \mapsto a$. It remains to show that
\[
  \dim \Ima(di)_D = \dim \Ker(df^{(r)})_D.
\]

Identify $T_a(J)^{\vee}$ with $\Gamma(C, \Omega^1)$ using the isomorphisms
arising from (2.1). Then (5.3) shows that $\omega$ is zero on the image of
$T_D(C^{(r)})$ if and only if $(\omega) \ge D$, that is,
$\omega \in \Gamma(C, \Omega^1(-D))$. Therefore the image of $(df^{(r)})_D$
has dimension $g - h^0(\Omega^1(-D)) = g - h^1(D)$, and so its kernel has
dimension $r - g + h^1(D)$. On the other hand, the image of $(di)_D$ has
dimension $|D|$. % [sic: printed "dimension |D|", meaning the dimension of the linear system |D|]
The Riemann-Roch theorem says precisely that these two numbers are equal, and
so completes the proof.

\paragraph{Corollary 5.5.}
\textit{For all $r \le g$, $f^r \colon C^r \to W^r$ has degree $r!$.}

\paragraph{Proof.}
It is the composite of $\pi \colon C^r \to C^{(r)}$ and $f^{(r)}$.
\hfill $\square$

\paragraph{Remark 5.6.}
(a) The theorem shows that $J$ is the unique abelian variety birationally
equivalent to $C^{(g)}$. This observation is the basis of Weil's construction
of the Jacobian. (See \S 7.)

(b) The exact sequence in (5.1b) can be regarded as a geometric statement of
the Riemann-Roch theorem (see especially the end of the proof). In fact it is
possible to prove the Riemann-Roch theorem this way (see Mattuck and Mayer
1963).

(c) As we observed above, the fibre of
$f^{(r)} \colon C^{(r)}(k) \to J(k)$ containing $D$ can be identified with
the linear system $|D|$. More precisely, the fibre of the map of functors
$C^{(r)} \to J$ is the functor $\Div_C^D$ of (3.14); therefore the fibre of
$f^{(r)}$ containing $D$ (in the sense of algebraic spaces) is a copy of
projective space of dimension $h^0(D) - 1$. Corollary 3.10 of Chapter I shows
that conversely every copy of projective space in $C^{(r)}$ is contained in
some fibre of $f^{(r)}$. Consequently, the closed points of the Jacobian can
be identified with the set of maximal subvarieties of $C^{(r)}$ isomorphic to
projective space.

Note that for $r > 2g - 2$, $|D|$ has dimension $r - g$, and so
$(df^{(r)})_D$ is surjective, for all $D$. Therefore $f^{(r)}$ is smooth (see
Hartshorne 1977, III 10.4), and the fibres of $f^{(r)}$ are precisely the
copies of $\mathbb{P}^{r-g}$ contained in $C^{(r)}$. This last observation is
the starting point of Chow's construction of the Jacobian Chow 1954.
% [sic: printed as "the Jacobian Chow 1954." — hyperlinked citation without parentheses]
